
\newcommand{\astar}{a_*}
\newcommand{\abar}{\bar a}
\newcommand{\epsES}{\varepsilon_{\rm ES}}
\newcommand{\qhat}{\widehat q}
\newcommand{\calZ}{\mathcal Z}
\newcommand{\calS}{\mathcal S}

\section{13-state First-Visit MCES example}
\label{app:fv}

\subsection{MDP, cycle, and notation}

This appendix describes the discounted 13-state example with 
exploring starts.  The discount and exploring-start perturbation level are
\begin{equation*}
    \gamma=0.9999,
    \qquad
    \epsES=10^{-4}.
\end{equation*}
The "move" action is denoted by $\astar$ and the "stay" action by $\abar$.
A successful "move" sends
\begin{equation*}
    s_i\longmapsto s_{i+1},
\end{equation*}
with indices understood modulo $13$.  See Figure~\ref{fig: long MDP}.
Action $\astar$ gives reward $1$ and succeeds with probability $\tau$; on success it moves to the next state, and on failure the episode terminates.  Action $\abar$ gives reward $0$ and remains in the same state.

A deterministic greedy policy is represented by its "stay" set $S\subseteq\{s_1,\ldots,s_{13}\}$.  Its value function satisfies
\begin{equation*}
V_S(s_i)=
\begin{cases}
\gamma V_S(s_i), & s_i\in S,\\[1mm]
1+\gamma\tau V_S(s_{i+1}), & s_i\notin S,
\end{cases}
\end{equation*}
and the exact policy action values are
\begin{equation*}
q_S(s_i,\astar)=1+\gamma\tau V_S(s_{i+1}),
\qquad
q_S(s_i,\abar)=\gamma V_S(s_i).
\end{equation*}
Because $\gamma<1$, "stay" states satisfy $V_S(s_i)=0$.

The intended adjacent cycle is
\begin{equation}
\label{eq:intended-cycle}
\{s_1\},\{s_1,s_{13}\},\{s_{13}\},\{s_{13},s_{12}\},\ldots,
\{s_2\},\{s_2,s_1\}.
\end{equation}
Equivalently, the local pattern is
\begin{equation*}
    \{s_i\}\longrightarrow \{s_i,s_{i-1}\}\longrightarrow \{s_{i-1}\}.
\end{equation*}
The normal start laws will be chosen later.

There are two requirements for a deterministic cycle. Firstly it must
converge to a point on the boundary and secondly it must follow the
right set of greedy policies as it does so. The first condition we
call the full cycle balance condition and the second the sign-margin
criterion. These are the subjects of the next two sections.

\subsection{General full-cycle balance equations}
\label{sec:balance-general}
We require that, provided the trajectory follows the right policies,
it converges onto the boundary.

Let $\calZ=\{(s_i,a):1\le i\le13,\ a\in\{\astar,\abar\}\}$ be the set
of state-action starts.  For phase $j$ in the intended cycle, let
$S_j$ be the "stay" set, $q_j=q_{S_j}$ its exact action-values, and
$\nu_j$ its nominal start distribution on $\calZ$. We impose a
rotational symmetry, that 
$\nu_{\{s_{i}\}}(s_{i+k},a)$ is a function of $k$ only, so it is only
necessary to specify one distribution over $\calZ$.
With
exploring-starts the  mixture is
\begin{equation*}
    \nu_j^{\epsES}=(1-\epsES)\nu_j+\epsES U_{\calZ},
    \qquad
    U_{\calZ}(z)={1\over 26}.
\end{equation*}
Let $\mu_j^z(s,a)$ be the probability of visiting $(s,a)$ when the episode starts from $z$ and follows policy $S_j$ after the start action.  The phase-$j$ first-visit probability under the perturbed start law is
\begin{equation*}
    \mu_j(s,a)=\sum_{z\in\calZ}\nu_j^{\epsES}(z)\mu_j^z(s,a).
\end{equation*}

For one complete pass through the 26 intended phases, define
\begin{equation*}
    B_a(s)=\sum_{j=1}^{26}\mu_j(s,a),
    \qquad
    A_a(s)=\sum_{j=1}^{26}\mu_j(s,a)q_j(s,a).
\end{equation*}
Here $B_a(s)$ is the expected number of first-visit updates to $(s,a)$ per normal full cycle, and $A_a(s)$ is the corresponding expected return-weighted update mass.

If fractional expected updates are applied over one full cycle, then
\begin{equation*}
    N^+(s,a)=N(s,a)+B_a(s)
\end{equation*}
and
\begin{equation*}
    q^+(s,a)=
    {N(s,a)q(s,a)+A_a(s)\over N(s,a)+B_a(s)}.
\end{equation*}
Equivalently,
\begin{equation*}
    q^+(s,a)-q(s,a)
    ={B_a(s)\over N(s,a)+B_a(s)}
    \left({A_a(s)\over B_a(s)}-q(s,a)\right).
\end{equation*}
Thus the full-cycle target for coordinate $(s,a)$ is $A_a(s)/B_a(s)$.

The boundary cancellation condition asks that, for every state,
\begin{equation}
    {A_{\astar}(s)\over B_{\astar}(s)}=
    {A_{\abar}(s)\over B_{\abar}(s)}=
    \qhat.
\label{eq:abstract-balance}
\end{equation}
i.e. the target is on the boundary. By rotational symmetry it is enough to impose the scalar equation at one state,
\begin{equation}
    F:=A_{\astar}(s_1)B_{\abar}(s_1)-A_{\abar}(s_1)B_{\astar}(s_1)=0,
\label{eq:abstract-F}
\end{equation}
The numerator of $F$ is a high order polynomial in $\tau$ and a low
order polynomial in each component of $\nu$.

The corresponding count ratio is
\begin{equation}
    r={B_{\astar}(s)\over B_{\abar}(s)}.
\label{eq:abstract-ratio}
\end{equation}
Choosing initial counts in this ratio makes the two action coordinates have matching first-order full-cycle learning rates:
\begin{equation*}
    {B_{\astar}(s)\over N_0(s,\astar)}=
    {B_{\abar}(s)\over N_0(s,\abar)}.
\end{equation*}

\subsection{Phase-sign margin LP}
\label{sec:sign-margin-lp}

The balance equations center the deterministic cycle on the boundary, but they do not by themselves ensure that the intended greedy policies occur in the right order.  Let
\begin{equation*}
    \sigma_j(s)=
    \begin{cases}
    +1, & s\in S_j,\\
    -1, & s\notin S_j.
    \end{cases}
\end{equation*}
At the boundary level $\qhat$, define the scaled expected gap increment in phase $j$ by
\begin{equation}
    g_j(s)=
    \mu_j(s,\abar)\bigl(q_j(s,\abar)-\qhat\bigr)
    -{\mu_j(s,\astar)\over r}\bigl(q_j(s,\astar)-\qhat\bigr).
\label{eq:gap-increment}
\end{equation}
Let
\begin{equation*}
    C_k(s)=\sum_{j<k}g_j(s)
\end{equation*}
be the cumulative gap displacement before phase $k$.  The scaled
initial gap $d_0(s)$ is then chosen by
\begin{equation}
\begin{aligned}
\text{maximize}\quad & \rho,\\
\text{over}\quad & d_0\in\mathbb R^{13},\ \rho\in\mathbb R,\\
\text{subject to}\quad
&\sigma_k(s)\bigl(d_0(s)+C_k(s)\bigr)\ge \rho,
\qquad k=1,\ldots,26,
\quad s=s_1,\ldots,s_{13}.
\end{aligned}
\label{eq:sign-lp}
\end{equation}
A positive optimal value $\rho$ means that the phases are followed in
the intended order.

The initial points are then placed symmetrically around the boundary:
\begin{equation*}
    q_0(s,\astar)=\qhat-{d_0(s)\over 2N_0(s,\abar)},
    \qquad
    q_0(s,\abar)=\qhat+{d_0(s)\over 2N_0(s,\abar)}.
  \end{equation*}
The factor $1/2$ splits the desired total action gap equally above and below $\widehat q$.



\subsection{Rotational start search and selected start law}
\label{sec:start-search}
The initial search for a suitable start distribution $\nu$ was over a
wide class of rotationally symmetric distributions. Distributions with
small support were preferred, as these would avoid adding too much
extra noise into the full simulations. We started by
considering all $13^2$ pairs of deterministic starts. For each, we solved
\eqref{eq:abstract-F} for $\tau$, and for any roots in $(0,1]$ we
then solved  \eqref{eq:sign-lp} for $\rho$. There were no solutions
with a positive $\rho$.

We then broadened the search to distributions with a two point support
in the two-"stay" phases. That is, 
for offsets $o_1,o_2,o_3$ we set
\begin{align}
    \nu_{\{s_i\}}^{(o_1)}
    &=\delta_{(s_{i+o_1},\abar)},\\
    \nu_{\{s_i,s_{i-1}\}}^{(o_2,o_3,\lambda)}
    &=\lambda\delta_{(s_{i+o_2},\abar)}+(1-
    \lambda)\delta_{(s_{i+o_3},\astar)}.
\end{align}
For a fixed triple and trial $\tau$, equation~\eqref{eq:abstract-F} is
solved for a root $\lambda\in[0,1]$ and then the sign-margin
LP~\eqref{eq:sign-lp} is evaluated. The $\tau$ which gives the maximum margin $\rho$ is then
found by gridding. The offsets
$    (o_1,o_2,o_3)=(-6,+1,-1)$ were chosen as giving a good margin
with the corresponding $\lambda$ close to 1.

Thus
\begin{align}
\nu_{\{s_i\}}&=\delta_{(s_{i-6},\bar a)},\\
\nu_{\{s_i,s_{i-1}\}}&=\lambda\delta_{(s_{i+1},\bar a)}+(1-
\lambda)\delta_{(s_{i-1},a_*)}.
\end{align}
and the corresponding $\tau$ and $\lambda$ are given by
\begin{equation*}
    \tau=0.512333524979486,
    \qquad
    \lambda=0.923714016854384.
\end{equation*}
The resulting boundary quantities are
\begin{equation*}
    \qhat=1.971056244749326
    \qquad
    q_*={1\over 1-
    \gamma\tau}=2.050366396036596
    \qquad
    r=2.010039006519156.
\end{equation*}
The sign-margin LP gives
\begin{equation*}
    \rho=8.527647265071263\times10^{-3}
  \end{equation*}
There were no  distributions with a two point support
in the single-"stay" phases and a one point support in the two-"stay" phases which admitted a positive margin.

Note that this procedure is not specific to the 13 state MDP, it in fact gives a suitable three point start distribution for the corresponding cycle on examples with 5 states and above, although a larger number of states gives a larger margin.

Figure~\ref{fig:spiral} in the main text shows the resulting
deterministic cycle. Note that the corners of the cycle contract
linearly towards $(\widehat q,\widehat q)$, and so stay in the right greedy region, thus the cycles will persist indefinitely. The next section confirms this contraction.
 

\subsection{Full-cycle linear interpolation toward $\widehat q$}
\label{sec:linear-interpolation}

Note that the full-cycle balance equations also imply a geometric contraction.  Suppose counts are exactly proportional to the full-cycle masses,
\begin{equation*}
    N_0(s,\astar)=c_sB_{\astar}(s),
    \qquad
    N_0(s,\abar)=c_sB_{\abar}(s).
\end{equation*}
Then after $m$ complete normal cycles,
\begin{equation*}
    N_{26m}(s,a)=(c_s+m)B_a(s).
\end{equation*}
Using $A_a(s)=\qhat B_a(s)$, the next full-cycle endpoint is
\begin{align}
q_{26(m+1)}(s,a)
&={ (c_s+m)B_a(s)q_{26m}(s,a)+A_a(s) \over (c_s+m+1)B_a(s)}\\
&={c_s+m\over c_s+m+1}q_{26m}(s,a)+{1\over c_s+m+1}\qhat.
\end{align}
Therefore the two-action point
\begin{equation*}
    P_m(s)=\bigl(q_{26m}(s,\astar),q_{26m}(s,\abar)\bigr)
\end{equation*}
satisfies
\begin{equation}
    P_{m+1}(s)=(1-\alpha_m)P_m(s)+\alpha_m(\qhat,\qhat),
    \qquad
    \alpha_m={1\over c_s+m+1}.
\label{eq:linear-interpolation}
\end{equation}
Equivalently,
\begin{equation*}
    P_m(s)-(\qhat,\qhat)=
    {c_s\over c_s+m}\bigl(P_0(s)-(\qhat,\qhat)\bigr).
\end{equation*}
Thus the full-cycle endpoints move on the line segment toward the boundary point. The intermediate phase points need not lie precisely on this line, but deviate only by order $1/k$.
\subsection{First-visit probabilities}

Fix an anchor state $s_i$ and write $d=(m-i)\bmod 13$, so $s_m=s_{i+d}$.  Let $\mu_1^\varepsilon$ and $\mu_2^\varepsilon$ denote the perturbed first-visit probabilities in the one-"stay" and two-"stay" phases for the selected start law. The values for the current $\tau$ and $\lambda$ are given in the following table for convenience. 
\begin{table}[H]
\centering
\small
\caption{Single-phase first-visit probabilities by relative distance $d$ for the tuned $n=13$ start law.}
\vspace{0.3cm}
\label{tab:first-visit-probabilities}
%\resizebox{\textwidth}{!}{%
\pgfplotstabletypeset[
  col sep=comma,
  columns/d/.style={int detect,column name={$d$}},
  columns/{mu1_move}/.style={sci,sci zerofill,precision=6,column name={$\mu_1^\varepsilon(d,\textup{"move"})$}},
  columns/{mu1_stay}/.style={sci,sci zerofill,precision=6,column name={$\mu_1^\varepsilon(d,\textup{"stay"})$}},
  columns/{mu2_move}/.style={sci,sci zerofill,precision=6,column name={$\mu_2^\varepsilon(d,\textup{"move"})$}},
  columns/{mu2_stay}/.style={sci,sci zerofill,precision=6,column name={$\mu_2^\varepsilon(d,\textup{"stay"})$}},
  every head row/.style={before row=\toprule,after row=\midrule},
  every last row/.style={after row=\bottomrule}
]{figures/first-visit-probabilities.csv}
\end{table}







\subsection{Sampled recovery strategy}
\label{sec:recovery-strategy}

The balance equations and margin LP describe the intended normal
cycle. In the sampled first-visit MC run, however, return noise,
trajectory noise, and rare exploring starts can move the greedy policy off the intended 26-phase cycle. The simulation therefore has to choose start distributions even when the observed greedy "stay"-set is not one of the intended cycle phases.

Let
\begin{equation*}
    S=\{s:q(s,\bar a)>q(s,a_*)\}
\end{equation*}
denote the set of states for which the "stay" action is greedy at the start of an update. If $S$ is one of the intended cycle phases, the controller uses the normal start law from Section~\ref{sec:start-search}.
If $S$ is off-cycle, the rollout policy is the current greedy policy with "stay" set $S$, and the controller chooses from the $26$ state-action starts as follows:

Let the current action gap be
\begin{equation*}
    G(s)=q(s,\bar a)-q(s,a_*).
\end{equation*}
If $S$ is empty then all $G(s)$ are negative, in this case we choose $s_0=\arg\max_s G(s)$ as the state whose action values are nearest the "move"/"stay" boundary. For other non-cycle phases then for $s_i\in S$ we let $m(s_i)= {\mathbf 1}(s_{i-1}\in M)+ {\mathbf 1}(s_{i+1}\in M)$ denote the number of neighbours of $s_i$ in $S$ and the $S_{\text{min}}=\{s\in S:m(s)=\min m\}$. That is, $S_{\text{min}}$ is the subset of $S$ consisting of the states with the minimum number of neighbours (all operations are modulo 13 and so $s_{1}$ and $s_{13}$ are considered neighbours). Finally we let $s_0=\arg\min_{s\in S_{\text{min}}} G(s)$
be the "stay" state with the fewest number of "stay" neighbours whose action values are closest to the "move"/"stay" boundary. This guarantees, for example, that if $M$ consists of an adjacent pair and one or more isolated states then one of the isolated states is chosen, and in a run of adjacent states then the two ends are chosen in preference to an internal state. Finally, the starting pair is chosen as the "stay" action in this state $(s_0,\bar a)$, in an attempt to move the action value across the boundary. If successful this would create a "stay" state if $S$ is empty or reduce a "stay" state, moving $S$ closer to to a cycle phase, otherwise. If at the next step $S$ is still not a cycle phase then the process is repeated.      

For the simulations we present, approximately $83.6\%$ of updates are
normal on-cycle updates, $14.9\%$ are attempted recoveries from the
optimal all-"move" policy and the remaining $1.5\%$ are recoveries from
all other non-cycle policies.  

\textit{ChapGPT 5.5 Pro was used to help refine and help code this example. The core simulation was hand-coded.}
